\section{Framework and Estimands}
\label{sec:model}
\StageThreeSectionNote

This paper studies coordination persistence as a risk-set-defined continuation
problem on non-root comments in Moltbook. Each non-root parent comment creates a
follow-up opportunity: another comment either arrives as a direct reply within
an operational horizon or it does not. The central OR/MS question is therefore
not whether activity exists somewhere on the platform, but which margin limits
within-horizon continuation on a common opportunity set---reply incidence or
conditional reply speed. The framework below defines that control panel, links
it to thread-level coordination outcomes, and states the exact algebraic
propositions that discipline estimation and intervention-priority logic.

\subsection{Scientific object and scope}
\label{sec:model:scope}

Let \(m\) index a non-root comment that can itself receive a direct reply.
Write \(\tau_m\) for its creation time and \(T_m\) for the lag from \(\tau_m\)
to the first direct reply, if one occurs. Let \(C_m\) denote observed follow-up
from \(\tau_m\) until administrative end, archive interruption, or other
coverage loss, and define
\[
\delta_m := \mathbf{1}\{T_m \le C_m\}.
\]
All primary incidence and timing estimands are defined on these non-root parent
comments; direct replies to root posts are outside the primary persistence
object. This scope isolates continuation beyond the first layer and avoids
conflating root attraction with sustained comment-to-comment coordination.

The paper is observational throughout. It measures continuation and its
structural consequences in archived interaction traces, but it does not identify
causal intervention effects, latent exposure, or pure autonomous intent. These
boundaries are important because the same observed reply pattern can be
generated by multiple latent mechanisms, including resurfacing, visibility,
moderation, memory loss, or scheduled agent availability
\citep{hawkes1971spectra,whitt2004efficiency,reed2012hazard,jouini2010online}.

\subsection{Observed thread state and downstream coordination outcomes}
\label{sec:model:branching}

Let \(j \in \mathcal{J}\) index a thread with root post \(n=0\) and comments
\(n=1,\ldots,N_j\). Event \(n\) is \((t_{jn},a_{jn},p_{jn})\), where
\(t_{jn}\) is the UTC timestamp, \(a_{jn}\) is the observed author identifier,
and \(p_{jn}\in\{0,\ldots,n-1\}\) is the parent index. Node depth is defined
recursively by
\begin{equation}
\label{eq:depth}
d_{j0}=0,
\qquad
d_{jn}=d_{j,p_{jn}}+1 \quad (n\ge 1),
\end{equation}
and maximum thread depth is \(D_j:=\max_{0\le n\le N_j} d_{jn}\).

For node \((j,n)\), let
\[
c_{jn}:=\#\{r>n : p_{jr}=n\}
\]
be its direct-child count. The branching-by-depth profile is
\begin{equation}
\label{eq:branching-by-depth}
\bar c_k := \mathbb{E}\!\left[c_{jn}\mid d_{jn}=k\right], \qquad k\ge 0.
\end{equation}
This profile distinguishes root-heavy star patterns from deeper continuation.

Define the directed reply-edge set
\[
\begin{aligned}
E_j
&:=\{(u,v): \exists n\ge 1 \text{ with } a_{jn}=u,\\
&\hspace{2.7em} a_{j,p_{jn}}=v,\ u\neq v\},
\end{aligned}
\]
and let
\[
\begin{aligned}
\Delta_j
&:=\{\{u,v\} : (u,v)\in E_j\\
&\hspace{2.7em} \text{or } (v,u)\in E_j\}
\end{aligned}
\]
be the unordered dyads with at least one directional reply. Thread-level
reciprocity is
\begin{equation}
\label{eq:reciprocity-rate}
R_j
:=\frac{1}{|\Delta_j|}\sum_{\{u,v\}\in \Delta_j}
\mathbf{1}\{(u,v)\in E_j \text{ and } (v,u)\in E_j\},
\end{equation}
with \(R_j\) left undefined when \(|\Delta_j|=0\).

Thread-level re-entry is
\begin{equation}
\label{eq:reentry-rate}
RE_j
:=\frac{\#\{n\ge 1 : a_{jn}\in\{a_{j1},\ldots,a_{j,n-1}\}\}}{N_j},
\end{equation}
again defined on the non-root comment sequence. Together,
\[
G_j := \left(D_j,\{\bar c_k\}_{k\ge 0},R_j,RE_j\right)
\]
summarizes the downstream coordination structure that accumulates from
micro-level continuation successes and failures
\citep{harris1963theory,gomez2013structure,aragon2017thread,meital2024branch}.

\subsection{The horizon-throughput control panel}
\label{sec:model:estimands}

For horizon \(h>0\), define the horizon risk-set indicator
\[
R_h(m):=\mathbf{1}\{C_m\ge h \text{ or } T_m\le h\}.
\]
The paper's primary object is the common-risk-set control panel
\[
\mathcal{C}_h := (\pi_h,\phi_h,q_h), \qquad h\in\{5\text{ min},1\text{ h}\},
\]
where
\begin{align}
\pi_h &:= \Pr(\delta_m=1 \mid R_h(m)=1), \\
\phi_h &:= \Pr(T_m\le h \mid \delta_m=1,\ R_h(m)=1), \\
q_h &:= \Pr(T_m\le h \mid R_h(m)=1).
\end{align}
Earlier drafts reported \(q_h\) under the notation \(p_h\); throughout the
final architecture, \(q_h\) is the preferred notation because it makes the
incidence--timing decomposition explicit.

The secondary descriptive incidence margin is
\[
p_{\mathrm{obs}}:=\Pr(\delta_m=1),
\]
the coverage-conditional ever-reply share. It is reported for descriptive
orientation only and is not the flagship decision object because it is sensitive
to heterogeneous follow-up.

Conditional timing is summarized nonparametrically through the distribution of
\(T_m\mid (\delta_m=1)\), especially
\((t_{50},t_{90},t_{95})\) and short-window masses such as
\(\Pr(T_m\le 30\text{ s}\mid \delta_m=1)\) and
\(\Pr(T_m\le 5\text{ min}\mid \delta_m=1)\).

\begin{proposition}[Common-risk-set factorization]
\label{prop:risk-set-factorization}
For any horizon \(h\) with \(\Pr(R_h(m)=1)>0\),
\begin{equation}
\label{eq:horizon-throughput-factorization}
q_h = \pi_h \phi_h.
\end{equation}
\end{proposition}

\begin{proof}
See Supplementary Material, Proposition 1.
\end{proof}

\subsection{Minimal mechanism scaffold and identification boundary}
\label{sec:model:mechanism}

The paper retains only a minimal observational mechanism scaffold. For parent
comment \(m\) and age \(u>0\), write
\begin{equation}
\label{eq:minimal-mechanism}
\lambda_m(u)=E_m(u)\,S_m(u),
\end{equation}
where \(E_m(u)\) is effective exposure or availability and \(S_m(u)\) is a
weakly decreasing staleness kernel. Heartbeat routines, resurfacing, ranking,
and platform visibility all load into \(E_m(u)\); memory decay, context loss,
and reply urgency load into \(S_m(u)\). This factorization is deliberately
agnostic: the current archives do not separately identify exposure from
staleness, so \Cref{eq:minimal-mechanism} is used as disciplined mechanism
vocabulary rather than as a fully estimated structural model
\citep{hawkes1971spectra,crane2008robust,cox1972regression}.

A parametric exponential kernel \(S_m(u)=\exp(-\beta u)\) is therefore only a
secondary timing diagnostic. When such a diagnostic is reported, the quantity
\(\ln 2/\beta\) is interpreted as an exponential-equivalent reply-hazard
timescale, not as a thread ``half-life'' and not as a primary estimand.

\subsection{Design assumptions and primary hypotheses}
\label{sec:model:assumptions}

The framework rests on the following design assumptions.

\begin{description}
\item[A1. Archive integrity and non-root scope.]
After deterministic deduplication and UTC normalization, parent links and event
ordering are accurate enough for aggregate tree reconstruction, and the primary
persistence object is restricted to non-root comments.

\item[A2. Risk-set identification and gap discipline.]
For fixed horizon \(h\), the quantities \(q_h\), \(\pi_h\), and \(\phi_h\) are
identified from parents with at least \(h\) observed follow-up or an observed
reply before \(h\). Coverage gaps are treated as missing exposure windows rather
than filled by imputation.

\item[A3. Cross-archive comparability.]
A larger public Moltbook archive can be restricted to a harmonized minimal
schema---thread identifier, comment identifier, parent identifier, author
identifier, UTC timestamp, and community label---sufficient for separate-archive
replication of the supported core metrics.

\item[A4. Mechanism non-separability.]
Without impressions, ranking, or exposure logs, visibility/availability effects
and staleness/memory effects are not separately identified from archival traces.

\item[A5. Local comparative statics.]
The intervention logic concerns local or one-margin perturbations to the
observed margins \((\pi_h,\phi_h)\); the archival data do not identify the
causal intervention effects that would generate those perturbations.

\item[A6. Coordination bridge.]
Higher sustained non-root horizon throughput should weakly support deeper
continuation, more non-root branching, more reciprocity, and more re-entry, but
the exact branching law is not identified.

\item[A7. Coarse heterogeneity proxies.]
Deterministic submolt labels are adequate for prespecified coarse heterogeneity,
whereas account-status fields are descriptive only and do not support mechanism
attribution.

\item[A8. Observed-dynamics and periodicity boundary.]
The paper makes claims about observed platform interaction dynamics only.
Aggregate periodicity diagnostics can detect strong synchronization when it is
present, but weak concentration does not rule out individual routines under
dephasing or jitter.
\end{description}

The empirical design tests the following hypotheses.

\begin{description}
\item[H1 (core).]
After gap-aware and horizon-standardized estimation, Moltbook exhibits a stable
low-incidence / fast-conditional regime on non-root reply opportunities at
\(h=5\) minutes and \(h=1\) hour. The same ordering is expected to survive
separate-archive replication on a larger public Moltbook archive once the
required minimal schema is verified.

\item[H2 (core).]
At minute-to-hour horizons, Moltbook is incidence-constrained rather than
timing-constrained: \(\phi_h\) materially exceeds \(\pi_h\), so equal-size or
moderate-cost-ratio perturbations imply larger local throughput gains from
incidence lift than from timing lift.

\item[H3 (core).]
Windows or strata with lower \(q_h\) also exhibit shallower depth tails, lower
non-root branching, lower reciprocity, and lower re-entry, consistent with
coordination-persistence limits.

\item[H4 (secondary).]
Prespecified coarse topic strata differ materially in \(q_h\) and in upper-tail
conditional timing, and the main direction of those contrasts is robust to
deterministic keyword-remapping sensitivity.

\item[H5 (diagnostic only).]
Aggregate four-hour periodicity is weak and non-load-bearing: departures from
exact phase uniformity, if detected, are too small to serve as the main
explanation of the persistence bottleneck.
\end{description}

\subsection{Operational decision rule}
\label{sec:model:or-diagnostic}

Because \Cref{eq:horizon-throughput-factorization} is exact, local one-margin
perturbations yield exact first-step gains. Let \(\Delta q_h^{(\pi)}\) denote
the change in \(q_h\) when only \(\pi_h\) changes by \(\Delta \pi_h\), holding
\(\phi_h\) fixed, and define \(\Delta q_h^{(\phi)}\) analogously when only
\(\phi_h\) changes.

\begin{proposition}[One-margin gains and cost-normalized priority]
\label{prop:local-gain-priority}
If only incidence changes, then
\[
\Delta q_h^{(\pi)}=\phi_h\,\Delta \pi_h.
\]
If only conditional timing changes, then
\[
\Delta q_h^{(\phi)}=\pi_h\,\Delta \phi_h.
\]
For positive intervention costs \(c_\pi\) and \(c_\phi\), the corresponding
cost-normalized local gains are
\begin{equation}
\label{eq:horizon-throughput-delta}
I_\pi:=\frac{\phi_h\,\Delta \pi_h}{c_\pi},
\qquad
I_\phi:=\frac{\pi_h\,\Delta \phi_h}{c_\phi},
\end{equation}
so local priority favors incidence lift if and only if \(I_\pi>I_\phi\).
\end{proposition}

\begin{proof}
See Supplementary Material, Proposition 2.
\end{proof}

\begin{remark}[Auxiliary coordination bridge]
The paper tests the link from \(q_h\) to downstream thread structure through
joint readouts of \((q_h,D_j,\bar c_k,R_j,RE_j)\). No identified branching
ratio or exact reproduction law is claimed. Any approximation such as
\(\Pr(D_j\ge K)\approx q_h^{K-1}\) is therefore heuristic and explicitly
non-load-bearing.
\end{remark}
